\subsection{Геометрия траекторий: окружность}

Для полёта по сложным траекториям часто используется параметрическое задание координат. Самый простой пример — окружность.

Окружность радиуса \(r\) в плоскости \(X\)–\(Y\) задаётся формулами:
\[
  x(\theta) = r \cdot \cos(\theta),\quad y(\theta) = r \cdot \sin(\theta),\quad z = \text{const}.
\]
Здесь \(\theta\) (тета) — текущий угол в радианах.

\begin{infobox}[Градусы и Радианы]
В математике и программировании тригонометрические функции (\texttt{sin}, \texttt{cos}) обычно принимают аргумент в \textbf{радианах}.
Если вы привыкли мыслить градусами (0\dots360), используйте формулу перевода:
\[
  \theta_{\mathrm{rad}} = \theta_{\mathrm{deg}} \cdot \frac{\pi}{180}
\]
Или встроенную функцию Lua: \texttt{math.rad(degrees)}.
\end{infobox}

\begin{codebox}[Реализация на Lua]
\begin{minted}{lua}
local angle = 0          -- угол в градусах
local r = 1.0            -- радиус 1 метр

-- Внутри таймера:
angle = angle + 5        -- шаг 5 градусов
local rad = math.rad(angle)
local x = r * math.cos(rad)
local y = r * math.sin(rad)
ap.goToLocalPoint(x, y, 1.0)
\end{minted}
\end{codebox}
